\chapter{Curriculum Vitae}
    \section*{Personal Information}
        \begin{tabular}{ll}
        Name: & John Doe \\
        Date of Birth: & 01.01.1990 \\
        Place of Birth: & Springfield, USA \\
        Nationality: & German \\
        \end{tabular}
    %
    \section*{Languages}
    \begin{tabular}{ll}
        German: & Native \\
        
        English: & (CEFR C1) Fluent \\
        
        Spanish: & (CEFR C2) Fluent \\
        
        French: & (CEFR A2) Basic \\
    \end{tabular}

    \section*{Distinctions and Awards}
        \begin{tabular}{lp{10cm}}
        %
        [2020] & Award for best M.Sc.~thesis in STEM disciplines, ``Unifreunde'' Association Bremen \\
                
        [2013] & Award for best high school diploma in the regency, Some other Association \\
        
        \end{tabular}
     \vspace*{\fill}
    
    \section*{Education}
        \begin{tabular}{lp{10cm}}
        [2021--ongoing] & \textbf{PhD in Electronic Structure Theory} \\
                     & Universität Bremen, Germany \\
                     & Supervisor: Lucio Colombi Ciacchi \\ \hline
        
        [2018--2020] & \textbf{M.Sc.~XXX} \\
                     & XXX University, Germany \\
                     & Thesis: ``Thesis title'' \\
                     & Supervisor: XYZ \\
                     & Grade: XXX \\ \hline
        
        [2015--2018] & \textbf{B.Sc.~XXX} \\
                     & XXX University, Germany \\
                     & Thesis: ``Thesis title'' \\
                     & Supervisor: XYZ \\
                     & Grade: XXX \\ \hline
                     
        [2007--2015] & \textbf{University Entrance Diploma (Abitur)} \\
                     & XXX High school, Germany \\
                     & Grade: 1.0 (full marks) \\ 
        \end{tabular}
        %
        
    \section*{Research Stays}
        \begin{tabular}{lp{10cm}}
        [09/2020]           & \textbf{Institution}: Host Institute A (Lausanne, Switzerland) \\
                            & \textbf{Host}: Host's Name  \\
                            & \textbf{Topic}: Short description \\ \hline
                            
        [08/2019]    & \textbf{Institution}: Host Institute B, (London, UK) \\
                            & \textbf{Host}: Host's Name \\
                            & \textbf{Topic}:  Short description \\ \hline
        
        [08--09/2016]       & \textbf{Institution}: Host Institute C, (Bogotá, Colombia) \\
                            & \textbf{Host}: Host's Name \\
                            & \textbf{Topic}: Short description \\ 
        \end{tabular}

\chapter{Research activities}

\nonumsection{Publications}
\begin{etaremune}
    \item[\cite{mackeOrbitalResolved2024}] \textbf{E.~Macke}, I. Timrov, N. Marzari, L. Colombi Ciacchi, \textit{Orbital-Resolved DFT+U for Molecules and Solids}, J.~Chem.~Theory Comput.~20, 11, 4824–4843 (2024)
    %
    \item [\cite{bastoneroFirstprinciples2025}] L. Bastonero, C. Malica, \textbf{E. Macke}, M. Bercx, S. Huber, I. Timrov, N. Marzari, \textit{First-principles Hubbard parameters with automated and reproducible workflows}, npj Comput.~Mater 11, 183 (2025)
    %
    \item [\cite{balzarettiMultiscale2025}] F. Balzaretti, M. von Einem, L. Gerhards, \textbf{E. Macke}, T. Neudecker, L. Colombi Ciacchi, W. Dononelli, S. Köppen-Hannemann, \textit{A Multiscale Modeling Approach to Understand the Degradation of Glyphosate under Dark and Light Conditions}, J.~Phys.~Chem.~C 129, 31, 13910–13922 (2025)
    %
    \item[\cite{wardaGetting2025}] \textbf{E.~Macke\textsuperscript{\textdagger}}, K. Warda\textsuperscript{\textdagger}, I. Timrov, L. Colombi Ciacchi, P. Kowalski, J.~Chem.~Theory Comput.~22, 2, 1016–1029 (2026)\footnote{\textsuperscript{\textdagger} Contributed equally.}
\end{etaremune}

\nonumsection{Talks (as a speaker)}
\begin{etaremune}
    \item \textbf{(invited)} \textit{Hubbard manifolds in DFT+U: Are we getting them right?}, Hubbard Workshop, Gandia, Spain (22-26 September 2025)
    \item \textbf{(contributed)} \textit{Orbital-resolved DFT+U for accurate predictions of Prussian blue analogues}, APS March Meeting, Anaheim/CA, United States of America (16-21 March 2025)
    \item \textbf{(invited)} \textit{Optimizing structures in
    theoretical materials science}, IEEE EDS SUMMER SCHOOL, Morelia, Mexico (12-15 June 2023)
    \item \textbf{(contributed)} \textit{Accurate spin-state energies from pinpointed DFT+$U_m$}, Canadian Chemistry Conference and Exhibition ,Vancouver, Canada (4-8 June 2023)
    \item \textbf{(contributed)} \textit{Pinpointing Hubbard corrections to tackle inhomogeneous nl subshells}, DPG Spring Meeting, Regensburg, Germany (4-9 September 2022)
    \item \textbf{(contributed)} \textit{The DFT$+U_m$ method}, 33rd IUPAP Conference on Computational Physics, Austin/TX, United States of America [conference held virtually] (1-4 August 2022)
    \item \textbf{(contributed)} \textit{On Hubbard models (and how to eventually overcome DFT’s flaws)}, Hike\&Science event ``Free-energy landscapes of bio-hybrid materials systems'', Gosau, Austria (26 June - 1 July 2022)
\end{etaremune}

\nonumsection{Posters}
\begin{etaremune}
    \item \textit{Pinpointing Hubbard corrections to tackle inhomogenous nl subshells: The DFT+$U_m$ method}, MAPEX Early Career Researcher workshop, Bremen, Germany (26 April 2022)
    \item \textit{Ab-initio Hubbard-parameters for battery materials}, RTG-QM3~review, Bremen, Germany (23 February 2021)
\end{etaremune}


\nonumsection{Teaching}
\begin{etaremune}
    \item \textit{Deep learning for predicting and analyzing EV charging stations} [conceptualization, lecturing, supervision of hands-on sessions and evaluation], computer laboratory, 4.5 ECTS, summer term 2025
    \item \textit{Surfaces and Interfaces} [conceptualization and supervision of hands-on sessions], lecture and computer laboratory, 6 ECTS (with L. Colombi Ciacchi), summer term 2023
    \item \textit{From molecular dynamics to physical properties} [conceptualization and supervision of hands-on sessions], STEM lab course for high-school students (with P. Goetz and L. Colombi Ciacchi), 2-3 November 2023
    \item \textit{Efficient data processing and dynamical visualization using Python} [conceptualization, lecturing, supervision of hands-on sessions and evaluation], lecture and computer laboratory, 4.5 ECTS (with W. Dononelli), summer term 2022
\end{etaremune}
